% Restructured 2026-09-28: the withdrawn and negative results, moved verbatim
% from method_tests.tex so the main text leads with the claims.
\subsection{The encoder's advantage does not survive}
\label{subsec:q1}

\paragraph{Pre-registered question.} Does the encoder's win over the
best baseline survive the move from the shared-initialisation cohort to
independently initialised ones? The rule, fixed in advance: for each base
model, compare the encoder against the champion baseline selected on the
mean over the three independent cohorts, with a tie threshold of $0.005$
nats and a one-directional noise gate at $2\times\mathrm{SE}$ that may
only downgrade a win or loss to a tie, never promote.

\begin{table}[t]
\centering
\small
\caption{Pre-registered replication on independently initialised cohorts,
and the two controls applied to it. Entries are the paired difference $d$
in worst-task NLL excess (nats), mean over three cohorts; positive favours
the encoder. \emph{As published} is the registered primary arm.
\emph{Matched} corrects an inference-path mismatch: at $\deff$ realisation
the merged adapter has rank $64$ against the baselines' $16$, and the fused
fast path does not take mixed-rank adapters, so the encoder alone ran
vanilla \textsc{transformers}; the two paths differ by up to $0.0105$ nats
on identical inputs. The rank-$16$ realisation is matched on both counts
and needs no new measurement. $T = 3$ additionally drops the adapter that
never learned its task (\S\ref{subsec:scope}). Levels are not comparable
across $T$, so only the verdict transfers. The noise gate
($2\times\mathrm{SE}$, one-directional) downgraded nothing; SE on $d$ runs
$0.0002$ to $0.0012$.}
\label{tab:q1}
\begin{tabular}{lcccl}
\toprule
base & as published & matched & matched, $T=3$ & verdict \\
\midrule
Llama-3.1-8B & $+0.0335$ & $+0.0257$ & $+0.0033$ & wins $\to$ \textbf{ties} \\
Mistral-7B   & $-0.0055$ & $-0.0022$ & $-0.0012$ & loses $\to$ \textbf{ties} \\
Qwen2.5-7B   & $-0.0002$ & $-0.0004$ & $+0.0009$ & ties throughout \\
Yi-1.5-9B    & $-0.0150$ & $-0.0154$ & $-0.0053$ & loses throughout \\
\midrule
\multicolumn{4}{l}{\emph{count}} & 1/1/2 $\to$ 1/2/1 $\to$ \textbf{0/3/1} \\
\bottomrule
\end{tabular}
\end{table}

\paragraph{Outcome: it does not survive.} The registered arm gives one win,
one tie and two losses, and each control removes another piece: matching
the inference path turns Mistral's loss into a tie, and dropping the
untrained adapter turns Llama's win into a tie, leaving \textbf{no wins on
any base}. Both pre-specified robustness lines agree with the primary arm.
Re-selecting the champion inside each cohort, the variant deliberately
biased against our method, returns the same one/one/two; applying the
original single-cohort rule three times succeeds on \texttt{indep1} and
fails on \texttt{indep2} and \texttt{indep3}, so the earlier optimistic
reading was the best of three draws. Two further details: the rank-$64$ realisation disagrees at $T=3$, keeping Llama at
$+0.0081$, and Mistral's champion changes from TIES to TVQ.

The registered $T = 4$ arm stays primary. The direction is consistent
across all three columns, and the reading we take is that the encoder's
advantage was substantially a property of the degenerate cohort it was
developed on. On that cohort $\bar H$ is ill-conditioned by four
orders of magnitude (\S\ref{subsec:conditioning}), which is exactly what
the ridge term was introduced to handle, so the method was tuned on the one
regime where its distinguishing component has work to do.
\S\ref{subsec:ridge} makes that concrete: the ridge gain is $0.19$ to
$9.73$ nats on the shared cohort and $0.00$ to $0.10$ on the independent
one. What does not survive is the claim that the encoder is the better
method.

\subsection{A claim we withdraw}
\label{subsec:withdrawn}

The result in \S\ref{subsec:q1} raises an obvious and much more
interesting possibility: if method rankings shift when initialisation
changes, then published merging benchmarks, which as far as we can tell
mostly use shared-seed cohorts, might be measuring initialisation
geometry rather than method quality. We described this on first sight as
the strongest result available to us.

We pre-registered a control to test it, and the control did not support
the claim. The control asks whether the method ranking is stable
\emph{within} the independent regime, across the three cohorts that
differ only in initialisation draw. If rankings are unstable within a
single regime, then ranking changes across regimes cannot be attributed
to the regime. Two of four base models had a non-unanimous top-1 method
across the three cohorts, above the threshold fixed in advance, so by our
own rule the attribution fails and \textbf{the claim that published
merging benchmarks are confounded by shared-initialisation geometry is
withdrawn}.

We record one caveat, and we deliberately did not let it reverse the
verdict. The control as written counts a top-1 change as instability
regardless of margin, and both flips involve methods separated by far less
than the tie threshold: on Qwen the top two are $0.0140$ and $0.0141$, a
coin flip. A margin-aware version of the control would very likely give a
different answer. But that version was not what we pre-registered, and
rewriting a decision rule after seeing which way it fell is exactly the
practice the pre-registration exists to prevent. We stated the defect and
left the verdict, and registered a \emph{second} control rather than
editing the first.

\paragraph{A second control, reported but not used to reinstate anything.}
We registered a margin-aware version of the rule afterwards, in the same
document that governs \S\ref{subsec:regmean}: a base counts as unstable
only if its top-1 method changes \emph{and} the margin between the top two
exceeds the tie threshold in at least one cohort where the change occurs.
It needs no new cells. Run as written it returns \textbf{zero of four}
bases unstable against two of four for the margin-blind rule, because both
flips sit well inside the threshold: Mistral's margins are $0.0004$,
$0.0003$ and $0.0011$ nats and Qwen's $0.0001$, $0.0006$ and $0.0003$,
against $0.005$.

Read literally, that rule would reinstate the benchmark-confounding claim
as provisional. \textbf{We decline to reinstate it}, and the reason is
worth stating plainly because it cuts against us. This control was
registered after its data had been seen. The registration says so about
itself, and we could bound it in various ways, but no amount of disclosure
changes the fact that we would be resurrecting a claim we had already
withdrawn using a rule written once we knew how the first one fell. That
is the single thing pre-registration exists to prevent, and a paper making
this much of the practice should not spend its credibility on a claim it
does not need. The prevalence audit already carries the load of showing
this matters outside our own repository.

The margin-aware numbers stay in the paper as a measurement, because they
say something true and useful: the ranking flips that triggered the
withdrawal are all inside noise. That is a statement about the fragility of
single-draw rankings, not about published benchmarks.

So the claim we make is this, and no more: method rankings on this metric
move between cohorts that differ only in initialisation draw, the movement
is within noise once margin is accounted for, and the collapsed regime that
would make such rankings hard to interpret is common in released adapters.
The benchmark-confounding claim stays withdrawn. We have not measured a
published merging benchmark, and we do not claim one is wrong.

What survives independently of all of it is weaker and still worth saying:
single-seed method rankings on this metric are not reproducible, so merging
comparisons should report multiple initialisation draws.

\subsection{The predicted rate exponent is unresolved}
\label{subsec:w3}

Theorem~\ref{thm:general} predicts distortion decaying as $2^{-2R/\deff}$,
and we registered that as a slope in $[-2.4, -1.6]$ on at least three of
four bases. \textbf{The sweep cannot resolve it, and two earlier readings
of it are withdrawn.} A previously reported falsification rested on
asymptotes computed at a different ridge value, and the refitted sweep
turns out to have too little dynamic range to measure an exponent at all:
only two of the six quantiser widths sit above the metric's own noise
floor, one of which the registered rule excludes as clipping-dominated.
By the rule as written the outcome is \textsc{unresolved}, and this
removes the only positive evidence for the rate term anywhere in this
paper. Appendix~\ref{app:rate} gives the sweep, both faults and the rule
applied to the result.

One operational regularity does survive at the precision available: with
the ridge on, every base is within $0.005$ nats of its own asymptote from
three bits upward, so a practitioner quantising a merged adapter can stop
at three or four bits regardless of any exponent.

\subsection{DARE: a negative result with a mechanism, and its limits}
\label{subsec:dare}

DARE~\citep{yu2024dare} drops a random fraction of task-vector entries
and rescales the survivors by $1/(1-p)$. Composed with task arithmetic it
does essentially nothing in our setting: across a density sweep from
$0.05$ to $0.5$ on four base models, it beats plain task arithmetic in
three of twenty cells and all three are within $0.0005$ nats, a fifth of
their tie threshold.

Unlike most negative results this one has a mechanism that predicts its
own shape. The rescaling makes the masked task vector an unbiased
estimator of the unmasked one, $\E[D_{\mathrm{DARE}}] = D_{\mathrm{TA}}$.
We state the consequence carefully, because the obvious version of the
argument is not available here: negative log-likelihood of an LLM is
\emph{not} convex in the weights, so Jensen's inequality does not apply
globally and we do not claim
$\E[L(D_{\mathrm{DARE}})] \geq L(D_{\mathrm{TA}})$ as a theorem. What
does hold is the local statement. Expanding $L$ to second order about
$D_{\mathrm{TA}}$, the first-order term vanishes by unbiasedness and the
penalty is $\tfrac12\tr(H\Sigma) + o(\|\Sigma\|)$ with
$\Sigma \propto (1-p)/p$ and $H$ the local Hessian. Where $H \succeq 0$,
which is the generic situation near a trained solution, this is
non-negative and grows as density falls. That is what we observe,
on three of four base models. We verified our implementation against the
reference one and confirmed the operator reproduces the analytic signature
$\sqrt{(1-p)/p}$ to three decimals, so this is a property of the method
and not of our code.

\paragraph{The negative result does not generalise, and we tested that.}
The Jensen argument applies only to the task-arithmetic composition, since
TIES is biased and the mask changes which coordinates survive trimming. We
pre-registered the DARE-TIES composition test before implementing it. The
outcome is mixed: DARE composed with TIES \emph{helps}
decisively on Llama ($+0.0462$ nats, against a gate of $0.0036$) and hurts
decisively on the other three. By the rule fixed in advance this is a
mixed verdict and neither direction is claimed.

So the correct scope for the negative result is DARE composed with task
arithmetic, not DARE. We flag that Llama is now the odd model out in two
separate pre-registered tests, here and in Table~\ref{tab:q1}, without
proposing a mechanism for it; a base-level moderator that neither
pre-registration anticipated deserves a designed experiment rather than
another post-hoc reading.

